\scanpage{76}
\mspage{II. 24}
% archive sheet number: 75
% lines 1–3
Considérons \deleted{\unsure{l'image de $C$}} le morphisme canonique
\[C\longrightarrow X\]
composé de l'inclusion $C\hookrightarrow X_K$ et $X_K-X$.
% lines 4–9
Nous supposons que \deleted{\unsure{le}} l'image\inserted{$x$} de $\xi$ dans $X$ est contenue dans celle de $C$. Ce qui signifie $\dim\overline{\{x\}}\geq1$); \deleted{\unsure{il en résulte}} comme l'extension 3.1.1. ne dépend pas, à isom. près, du choix de $\xi$, cela est une hypothèse \unsure{anodine|admissible}. Soit $\xi_1$ un point géométrique de $C$ au dessus de $\xi$
% lines 10–13
\[(3.1.2)\]
\diagram{\begin{tikzcd}[ampersand replacement=\&]
\xi_1 \arrow[r,dash] \arrow[d,dash] \& \xi \arrow[d]\\
C \arrow[r,dash] \arrow[d] \& X \arrow[d,dash]\\
\operatorname{Spec}K \arrow[r,dash] \& S
\end{tikzcd}} $*$
et soit $\bar K$ la clôture sép. de $K=k(t)$ dans $\kappa(\xi_1)$.
% lines 14–18
\deleted{On en conclut un diagramme \inserted{commutatif} d'inclusions de corps}
\diagram{\begin{tikzcd}[ampersand replacement=\&]
\deleted{k} \arrow[r,dash] \arrow[d,dash] \& \deleted{K=k(t)} \arrow[dd,dash]\\
\deleted{\bar k} \arrow[dr] \& \\
\& \deleted{\bar K}
\end{tikzcd}}
% lines 19–22
\deleted{Considérons aussi le composé $K\bar k=\unsure{L|l}$ dans $\bar K$, qui n'est autre que le corps $\kappa(\bar t)$, où $\bar t$ est le point générique de $T_{\bar k}$.}
% lines 23–25
Le diagramme (3.1.2) donne un diagramme commutatif, à lignes exactes :
